\section{A polynomial-threshold obstruction}
The following elementary polynomial-threshold obstruction is due to
Aspnes, Beigel, Furst, and Rudich~\cite{ABFR}; we include its proof.
\begin{lemma}[Polynomial thresholds cannot approximate parity]
\label{lem:ptf-parity}
On a cube of dimension $k$, for $0\le s<k$, every degree-$s$ polynomial threshold, with a fixed common convention
for zero ties, has full-parity error at least
\begin{equation}\label{eq:ptf-parity}
 2^{-k}\sum_{j=0}^{\lfloor(k-s-1)/2\rfloor}\binom kj.
\end{equation}
This lower bound tends to $1/2$ whenever $s=o(\sqrt k)$.
\end{lemma}
\begin{proof}
A sufficiently small constant perturbation resolves any zero ties
without changing the specified threshold classifier. Thus suppose
$P$ has no zeros. Put $t=\lfloor(k-s-1)/2\rfloor$. If its error set
had fewer than $\sum_{j\le t}\binom kj$ points, a nonzero polynomial
$Q$ of degree at most $t$ could be chosen to vanish on every error.
The function $fPQ^2$ would then be nonnegative everywhere and positive
somewhere. But $PQ^2$, after multilinear reduction, has degree at most
$s+2t<k$, so its parity Fourier coefficient
$\mathbb E fPQ^2$ is zero, a contradiction. This proves
\eqref{eq:ptf-parity}. The last assertion follows from the central
binomial bound $\max_j2^{-k}\binom kj=O(k^{-1/2})$.
\end{proof}

